% optimization/alignment/40_initial_guesses.tex

\section{Construction of Initial Guesses}

\subsection{Motivation}

The performance of SQP methods strongly depends on the choice of the
initial parameter vector.

In railway applications, initial data is often available in the form of:
\begin{itemize}
\item measurement polylines,
\item existing alignments,
\item partially structured design data.
\end{itemize}

\subsection{Input Representation}

\begin{definition}[Input polyline]
Let
\[
P = (p_0,\dots,p_M), \quad p_i \in \R^2
\]
denote a measured or imported polyline.
\end{definition}

\subsection{Arc-Length Parameterization}

\begin{definition}
Define cumulative arc-length
\[
s_0 = 0,
\quad
s_{i+1} = s_i + \|p_{i+1} - p_i\|.
\]
\end{definition}

This defines a parameterization of the input data.

\subsection{Initial Curvature Estimation}

\begin{definition}[Discrete curvature]
For three consecutive points, define:
\[
\kappa_i
\approx
\frac{\theta_{i+1} - \theta_i}{\Delta s_i},
\]
where
\[
\theta_i = \arg(p_{i+1} - p_i).
\]
\end{definition}

This provides a first estimate of curvature from geometric data.

\subsection{Smoothing}

Raw curvature estimates are typically noisy.

\begin{definition}[Smoothing operator]
Apply a smoothing filter:
\[
\kappa_i^{(0)}
=
\sum_j w_{ij} \kappa_j.
\]
\end{definition}

Possible choices:
\begin{itemize}
\item moving average,
\item Gaussian filter,
\item low-pass filtering.
\end{itemize}

\subsection{Initial Cant Estimation}

Cant may be initialized based on equilibrium:
\[
\phi_i^{(0)}
=
\arcsin\left(\frac{v^2 \kappa_i^{(0)}}{g}\right).
\]

Alternatively, cant may be:
\begin{itemize}
\item set to zero,
\item taken from input data,
\item approximated by design rules.
\end{itemize}

\subsection{Projection to Discrete Grid}

\begin{definition}
Interpolate the values \(\kappa_i^{(0)},\phi_i^{(0)}\)
onto the optimization grid:
\[
s_0,\dots,s_N.
\]
\end{definition}

\subsection{Heuristic Segmentation}

The input alignment may be segmented into:
\begin{itemize}
\item straight sections,
\item circular arcs,
\item transition regions.
\end{itemize}

\begin{definition}[Segment detection]
A segment is classified based on curvature:
\[
\kappa_i \approx 0 \quad \Rightarrow \text{straight},
\]
\[
\kappa_i \approx \text{const} \quad \Rightarrow \text{arc}.
\]
\end{definition}

This structure can be used to improve the initial guess and is
compatible with railway alignment reconstruction approaches such as
those discussed by Kufver
\cite{Kufver1997MathematicalDescription,Kufver2000RailwayAlignment}.

\subsection{Final Initial Guess}

\begin{definition}[Initial parameter vector]
The initial guess is
\[
\mathbf{x}_0 =
(\kappa_0^{(0)},\dots,\kappa_N^{(0)},
\phi_0^{(0)},\dots,\phi_N^{(0)}).
\]
\end{definition}

\subsection{Interpretation}

The initial guess combines:
\begin{itemize}
\item measured geometry,
\item smoothing,
\item physical approximation.
\end{itemize}

\begin{summary}
A robust initial guess transforms raw geometric input into a structured
parameter vector suitable for optimization.

This step is essential for bridging real-world data and numerical
methods.
\end{summary}

